\documentclass[12pt, addpoints]{exam}

% --- 控制解答顯示 (由 Python GUI 動態注入) ---
[[ print_answers_config ]]

% --- 紙張設定：B4 ---
\usepackage[paperwidth=25.7cm, paperheight=36.4cm, top=2cm, bottom=2.2cm, left=2cm, right=2cm, columnsep=1cm]{geometry}

% --- 中文字體設定 ---
\usepackage{fontspec}
\usepackage{xeCJK}
\setCJKmainfont[AutoFakeBold=true]{AR PL UKai TW}

% --- 常用套件 ---
\usepackage{siunitx, float, multicol, xcolor, amsmath, amssymb, graphicx, tabularx, multirow, tikz, enumitem, titlesec, booktabs, pgfplots, ulem}
\pgfplotsset{compat=1.18}
\usetikzlibrary{arrows.meta, calc, patterns, positioning}

\raggedcolumns
\setlength{\columnseprule}{0.4pt}
\renewcommand{\choicelabel}{(\thechoice)}
\renewcommand{\choiceshook}{\setlength{\leftmargin}{25pt}\setlength{\itemsep}{0pt}}

\extrafootheight{2.0cm}
\footer{}{\iflastpage{\vfill \centerline{\Large\textbf{--- 試題結束 ---}}}{\vfill \centerline{\Large\textbf{（尚有試題，請翻頁繼續作答）}}}}{第 \thepage\ 頁，共 \numpages 頁}
\firstpageheader{}{}{第 \thepage\ 頁，共 \numpages 頁}
\runningheader{[[ year_code ]][[ exam_title ]]}{[[ grade ]][[ course ]]}{第 \thepage\ 頁，共 \numpages 頁}

\begin{document}

% --- 試卷標頭 ---
\begin{center}
    \renewcommand{\arraystretch}{1.5}
    \begin{tabularx}{0.98\textwidth}{|*{6}{>{\centering\arraybackslash}X|}}
        \hline
        \multicolumn{6}{|c|}{\fontsize{18pt}{20pt}\selectfont \textbf{[[ school_name ]]\ [[ semester ]]\ [[ exam_title ]]}}\\
        \hline\large
        年級 & 科目 & 範圍 & 班級 & 座號 & 姓名 \\
        \hline
        [[ grade ]] & [[ course ]] & [[ test_range ]] & & & \\
        \hline
    \end{tabularx}
\end{center}

\vspace{0.8em}
\noindent [[ full_instruction_content ]] % 這裡會注入完整的計分說明
\vspace{0.5em}
\hrule
\vspace{1em}

% --- 題目區塊 ---
\ifprintanswers
    \begin{questions}
        [[ questions_content ]]
    \end{questions}
\else
    \begin{multicols*}{2}
    \begin{questions}
        [[ questions_content ]]
    \end{questions}
    \end{multicols*}
\fi

\end{document}